\section{Hydrodynamic, Operational, and Economic Performance}
\subsection{From surface condition to ship performance}
Hull roughness and fouling change resistance, required power, fuel consumption, and emissions. Roughness models, coating hydrodynamic measurements, and operating-ship observations provide convergent engineering rationale, but the magnitude depends on vessel, speed, wetted area, fouling state, and weather \cite{speranza_modelling_2019,yeginbayeva_experimental_2018,wang_drag_2018,hakim_investigation_2019}.

Scaling a localized coupon or panel measurement to ship performance is classically done through Granville's similarity-law procedure, which shifts the smooth-wall boundary-layer velocity profile by a measured roughness function $\Delta U^+$ \cite{granville_1958,granville_1987}; the ITTC frictional-allowance line traces to Townsin's correlation work, which is valid for moderate roughness but not for the multiscale structures of heavy biofouling \cite{townsin_mosaad_1985,townsin_2003}, and Candries' multi-parameter topography results show that texture and form, not height alone, govern the transfer \cite{candries_2001_thesis,candries_atlar_2005}. Polynomial fits have been used to map qualitative fouling-rating scales (for example the US Navy $fr_{NSTM}$ rating) onto $k_s$ \cite{schultz2007roughness}, but a direct, universal geometry-to-hydrodynamics translation remains unavailable \cite{speranza_modelling_2019}: viscoelastic biofilms deform under shear and interact non-monotonically with near-wall turbulence \cite{snowdon_surface_2023}, so static fouling-state photographs or coupon coverage should not be converted into a ship-scale $k_s$ or $\Delta C_F$ without dynamic boundary-layer or resistance measurement.

\subsection{ISO 19030 and ITTC translation}
The ship-scale endpoint should be expressed through a recognized performance framework rather than by transferring a coupon drag result directly to fuel use. ISO 19030-1 and ISO 19030-2 define the logic for measuring changes in ship-specific hull and propeller performance, including corrected speed--power observations and the separation of environmental and operational influences \cite{iso19030_1_2016,iso19030_2_2016}. ITTC performance-prediction and speed--power-trial procedures provide the complementary resistance, propulsion, model--ship correlation, and trial-correction framework \cite{ittc_speed_power}. In practice, the coating study should supply a time-resolved surface condition and roughness record, while the ship-performance study should report the ISO/ITTC-corrected change in speed loss at constant power or power increase at constant speed.

For a first-order expectation, use the ITTC--1957 friction line as a reference rather than treating it as a fouling law,
\begin{equation}
C_F = \frac{0.075}{(\log_{10}Re-2)^2},
\end{equation}
and define the measured total-resistance change as $\Delta R_T/R_T$. At fixed displacement and speed, the corresponding propulsive-power change is approximately $\Delta P/P \simeq \Delta R_T/R_T$; at fixed delivered power, the equivalent speed loss must be obtained from the corrected speed--power curve. If roughness or slime primarily changes frictional resistance, a useful decomposition is
\begin{equation}
\frac{\Delta P}{P} \simeq \phi_F\frac{\Delta C_F}{C_F}+\frac{\Delta R_{\mathrm{res}}}{R_T}, \qquad \phi_F=\frac{R_F}{R_T},
\end{equation}
where $\phi_F$ is the vessel- and condition-specific friction-resistance share. This equation makes the scale limitation explicit: a percentage reduction in coating or slime friction is not the same percentage reduction in total ship power unless friction dominates the resistance budget.

Canonical roughness and slime work provides the physical calibration for this translation: roughness changes the near-wall flow, while biofilms can change both equivalent roughness and the viscoelastic character of the interface \cite{schultz2007roughness,hunsucker_biofilm_2018,snowdon_surface_2023}. The expected benefit of a coating should therefore be reported as a measured delta relative to a defined reference coating, with the reference's roughness, fouling state, Reynolds-number range, and cleaning history stated. For example, an illustrative 5\% reduction in the frictional component with $\phi_F=0.70$ implies approximately a 3.5\% reduction in total resistance only if the residual-resistance term is unchanged; it is not a universal ship-fuel prediction. The ISO/ITTC result should be accompanied by uncertainty and sensitivity to weather correction, speed, draft, propeller condition, and fouling distribution.

\subsection{Worked ship-scale example: illustrative, not a measured result}
The following calculation makes the translation explicit without presenting an invented vessel measurement as corpus evidence. Consider a representative handysize bulk carrier with $L=180$ m at 14 kn. Assume that coating weathering or light slime produces an equivalent roughness increment $\Delta k_s=50~\mu$m, and that the associated friction change has been measured or conservatively assigned as $\Delta C_F/C_F=12\%$. Let the friction-resistance fraction be $\phi_F=C_F/C_T=0.75$, within the assumed illustrative range $0.70$--$0.80$. Then
\begin{align}
\frac{\Delta P}{P} &\simeq \phi_F\frac{\Delta C_F}{C_F} \\
&\simeq 0.75\times0.12=0.090\quad (9.0\%).
\end{align}
If the corrected speed--power relationship is locally approximated by $P\propto V^3$, the constant-power speed change is
\begin{equation}
\frac{\Delta V}{V}\simeq-\frac{1}{3}\frac{\Delta P}{P}\simeq-3.0\%.
\end{equation}
With a simple illustrative propagation of $\pm1.5$ percentage points on the power penalty, the corresponding speed-loss interval is approximately $-2.5$--$-3.5\%$. This example does not provide an absolute fuel loss because engine load, propulsive efficiency, draft, weather, and operating hours are not specified. In practice, ISO 19030 corrections and ITTC trial/performance procedures must replace the assumed $\Delta C_F/C_F$, $\phi_F$, and cubic local law with measured vessel data.

\subsection{Quantitative expectations and validation ladder}
The practical expectation for a new nanocomposite or fouling-release coating is a bounded chain of measurements: (i) coupon or panel roughness and friction under a stated flow condition; (ii) pre- and post-immersion roughness, slime coverage, and adhesion; (iii) towing-tank or equivalent hydrodynamic comparison against the reference coating; and (iv) ISO 19030/ITTC-corrected ship data. Canonical studies can justify the direction and mechanism of a penalty, but the ship-scale benefit should be calculated from the vessel's measured resistance decomposition. Fuel or emissions savings should be reported only after converting corrected power or speed changes through the vessel's operating profile; a coating's laboratory settlement reduction alone is insufficient.

\subsection{Condition-specific quantitative anchors}
The corpus contains useful anchors, but they are not a pooled engineering range. In a groomed fouling-release panel study, biofilm thickness on groomed panels was reported as $33.6\pm14.8~\mu$m, while the ungroomed condition was significantly thicker; associated drag forces were $0.75\pm0.09$ N and $1.09\pm0.06$ N for groomed and ungroomed biofilms, respectively \cite{hunsucker_biofilm_2018}. Another study reports a literature-level biofilm shaft-power penalty of approximately $1$--$18\%$, which should be treated as a contextual range rather than a transferable prediction for a particular vessel \cite{snowdon_surface_2023}. These values support the direction of benefit from grooming, but they do not define a universal drag-versus-LoF curve.

For cleaning, the calibrated waterjet study found visually effective cleaning at approximately $80$--$530$ Pa wall shear over its tested conditions, with maximum design values around $0.5$ kPa and $0.07$ MPa stagnation pressure; later low-standoff events reached approximately $1.3$ kPa and $0.17$ MPa \cite{oliveira_ship_2020}. The latter values should be treated as upper test conditions, not recommended universal limits. For environmental release, one 120-day batch study reported total copper concentrations reaching $178.7\pm43.7$ mg/L in estuarine water and $33.1\pm7.3$ mg/L in seawater for a non-nano copper coating; these are vessel-free batch concentrations, not area-normalized field fluxes \cite{miller_nano_2020}. In a separate year-scale panel study, an average copper release rate of approximately $13~\mu$g~cm$^{-2}$~day$^{-1}$ was estimated over 364 days under that study's coating and cleaning conditions \cite{oliveira_ship_2020}. The difference in units and exposure configuration is precisely why release values should not be merged without a mass-transfer and boundary-condition model.

\subsection{Maintenance and life-cycle value}
A coating should be evaluated over its service pathway: application quality, immersion ageing, cleaning, repair, dry-docking, replacement, and disposal. Life-cycle cost work and biocide-free coating assessments show why material price alone is insufficient \cite{korican_life-cycle_2024,venettacci_environmental_2023,oliveira_novel_2022}. A modestly more expensive surface may be preferable if it retains function, reduces cleaning, and lowers operational penalty; the corpus does not support a universal economic ranking.

Condition-based scheduling can be formulated as a dynamic-programming or machine-learning problem using speed-through-water, delivered power, weather, fouling observations, and cleaning history, but no verified scheduling-model study or performance range is present in the evidence matrix. Bayesian decision support is likewise a suitable framework for balancing non-indigenous-species risk, coating release, cleaning cost, and excess CO$_2$, but the evidence base does not provide a validated posterior model. These methods should therefore be reported as proposed decision architectures until their training data, calibration, and out-of-sample performance are available.

For implementation, the LoF posterior should be calibrated before it is connected to a maintenance trigger. The relevant state is $P(\mathrm{LoF}\geq3\mid\mathrm{image})$, not the argmax class alone, and the trigger should compare expected drag-fuel cost with cleaning and brush-damage cost. Expected Calibration Error, reliability diagrams, subgroup performance under turbidity and lighting regimes, and false-positive versus false-negative costs should be monitored over time. An uncalibrated network can create asymmetric operational harm: false positives cause unnecessary coating wear and captured-waste burden, while false negatives allow fouling growth and hydrodynamic penalty to accumulate.

\section{Environmental Risk and Life-Cycle Sustainability}
\subsection{Release pathways}
Potential source terms include dissolved metals or biocides, detached nanoparticles, polymer fragments, lubricant loss, and cleaning wastewater. Metal-release studies, paint-particle toxicity studies, and hydroblasting effluent analyses indicate that antifouling benefit and environmental burden must be assessed together \cite{miller_nano_2020,kim_qualitative_2024,soon_characterization_2021}.

The local evidence supports a 120-day batch release record and an area-normalized year-scale copper estimate, but it does not provide a matched static-versus-sailing flux pair. Static flux, sailing flux, zinc flux under both conditions, microplastic particle counts, heavy-metal loads, and validated $1$--$5~\mu$m filtration cutoffs are therefore NR in this review. They should be measured using the fractionated, area-normalized protocol in Table~\ref{tab:release_protocol}, with particle-size distribution and filter recovery reported explicitly.

\subsection{Standards alignment for release and ecotoxicity testing}
The proposed release workflow is a standards-aligned application framework, not a new parallel standard. Table~\ref{tab:standards_mapping} maps each tier to recognized methods and distinguishes what is complementary from what must remain method-specific. Dissolved biocide release should follow the applicable ISO 15181 part; substitute ocean water should be prepared or justified against ASTM D1141; abrasion should use ASTM D4060 and/or the applicable ISO 7784 procedure with the wheel, load, cycles, and coating reported; and ecotoxicity modules should map to OECD TG 201, 202, 203, and, for treatment residuals, TG 209 \cite{iso15181_1_2007,astm_d1141,astm_d4060,iso7784_1_2016,oecd_tg201,oecd_tg202,oecd_tg203,oecd_tg209}. Shear-induced wear fractionation, colloidal/particulate separation, and matched aged-coating exposure are complementary additions needed to represent marine service; they do not supersede the recognized methods or permit whole-paint-flake measurements to be reported as dissolved-ion release.

\subsection{Regulatory compliance boundaries}
Testing standards establish how an endpoint is measured; they do not by themselves authorize a coating, a biocide, or an in-water cleaning discharge. The IMO Anti-fouling Systems Convention restricts harmful organotin compounds and controls cybutryne-containing anti-fouling systems, while ship flag, port-state, coating documentation, and current amendments determine the applicable compliance evidence \cite{imo_afs_convention}. In the European Union, the Biocidal Products Regulation (EU) No.~528/2012 governs active-substance approval and product authorization; a laboratory antifouling result cannot be treated as evidence of market authorization or unrestricted use \cite{eu_bpr_528_2012}.

In-water cleaning is additionally jurisdiction- and permit-dependent. Australian and New Zealand biofouling-management requirements and California State Lands Commission rules can impose vessel-specific documentation, approved methods, capture, treatment, and discharge conditions \cite{anz_biofouling_standards,california_slc_cleaning}. Capture targets such as at least 95\% particulate recovery in a specified size range, including jurisdictional ranges around 1--10~$\mu$m where applicable, must be verified against the current permit or rule rather than generalized across all ports. The engineering protocol is therefore complementary to regulation: it measures dissolved, colloidal, and particulate fractions and documents recovery efficiency, but the operator must still satisfy the controlling local standard, waste classification, biosecurity requirement, and discharge permit.

\subsection{Non-target effects and uncertainty}
The evidence base contains organism-specific toxicity and ecological observations, but endpoint comparability is limited. Terms such as non-toxic, green, or eco-friendly should therefore be reserved for assessments covering relevant exposure and life-cycle dimensions. Biocide-free systems avoid a deliberate biocide source but can still generate persistent polymers, nanomaterials, or cleaning-derived particles.

Carbon nanomaterial toxicology should be assessed over the coating lifecycle rather than only with pristine particles. Proposed mechanisms include reactive-oxygen-species generation, oxidative stress, lysosomal membrane damage, and altered autophagic processing, but the cited external toxicology record is not present in the available evidence base and no mechanism-specific chronic range is reported in the evidence matrix. These endpoints are therefore hypotheses for standardized chronic exposure testing, with dissolved-ion, particle, matrix-fragment, and aged-coating fractions separated.

\subsection{Asset-specific interpretation}
A hull prioritizes roughness, drag, cleaning access, and invasive-species control; an aquaculture net prioritizes flexibility, light, abrasion, and organism exposure; a membrane prioritizes flux and biofilm selectivity; and a heat exchanger prioritizes heat transfer and cleanability. Cross-asset claims require explicit justification.
